\documentclass[../document.tex]{subfiles}

\begin{document}
\par Результатом учебной практики стала модель оптического распознавания символов, способная решать тесты Captcha.
\par В основе модели лежит комбинация свёрточной и рекуррентной нейронных сетей, с использованием функции оценки потерь на основе коннекционистской темпоральной классификации.
\par Для реализации была использована библиотека TensorFlow, с API Keras на языке Python. Для обучения был использован набор данных, состоящий из 1040 Captcha-изображений.
\subsection{Описание изучаемой темы и области  применения}
\par В ходе учебной практики изучены методы распознавания цифровых изображений с помощью машинного обучения. Компьютерное зрение применяется во множестве областей, начиная оцифровкой книг и заканчивая самоуправляемыми автомобилями и медициной.
\par Компьютерное зрение также можно использовать для решения Captcha, специальных тестов Тьюринга, используемых для определения того, является ли пользователь человеком или компьютером, чтобы предотвратить атаки и спам. Долгое время Captcha состояли из искажённого текста, который был трудным для алгоритмов оптического распознавания символов тех времён. Но с использованием современных методов машинного обучения такие тесты можно легко решить, из-за чего современные Captcha больше не используют текст для своих тестов, и вместо него берут сложные изображения. Тем не менее, распознавание текстовых тестов Captcha служит полезным примером использования нейронных сетей для решения сложных задач.
\subsection{Изучение основных приёмов работы с машинным обучением}
\par Было изучено множество методов реализации и разновидностей нейронных сетей, их особенности и области применения. В ходе планирования была составлена следующая структура модели нейронной сети:
\par На вход нейронная сеть получает заранее обработанное изображение в виде матрицы со значениями от 0 до 1, где 0 - чёрный цвет, а 1 - белый.
\par На первом этапе используются слои свёрточной нейронной сети (Convolutional neural network, CNN) для уменьшения объёма данных и выделения особенностей и деталей на изображении. Этот этап позволяет следующим слоям работать с более абстрактной репрезентацией изображения, содержащей только представляющие интерес особенности.
\par На втором этапе слои рекуррентной нейронной сети (Recurrent neural network, RNN) испольщуются для нахождения закономерностей в данных, полученных свёрточными слоями и генерации потенциальных расшифровок изображения. Рекурретные слои будут считывать изображение слева направо, используя ширину в качестве измерения времени.
\par Для оценки выходных данных рекуррентной сети используется коннекционистская темпоральная классификация (Connectionist temporal classification, CTC). Это функция оценки потерь, адаптированная для данных, меняющихся со временем (в качестве которого служит ширина изображения), отражающая степень сходства выхода нейронной сети и предсказываемой строки.
\subsubsection{Практическая работа}
\par Сперва был импортирован набор данных с ресурса Kaggle (Приложение 1.2). Создан документ Jupyter Notebook, в котором будет содержаться программный код на языке Python, пояснения к коду и выходные данные.
\\
\begin{figure}[H]
  \centering
  \includegraphics[width=\textwidth]{dataset.png}
  \caption{Пример вводных данных}
\end{figure}
\par Реализация программы состояла из нескольких этапов:
\begin{enumerate}
  \item Предварительная обработка изображений.
\par На этом этапе считывается весь набор данных, состоящий из 1040 изображений в формате PNG. Устанавливается размер партии для обучения и валидации, равный 16, размер субдескритизации, равный 4, и максимальная длина строки с ответом на Captcha, равная 5.
\begin{minted}[fontsize=\footnotesize,breaklines]{python}
# Путь к папке с данными
data_dir = Path("./captcha_images_v2/")

# Список всех картинок
images = sorted(list(map(str, list(data_dir.glob("*.png")))))
labels = [img.split(os.path.sep)[-1].split(".png")[0] for img in images]
characters = set(char for label in labels for char in label)
characters = sorted(list(characters))

print("Картинок найдено: ", len(images))
print("Слов найдено: ", len(labels))
print("Уникальных символов: ", len(characters))
print("Имеющиеся символы: ", characters)

# Размер партии для обучения и валидации
batch_size = 16

# Требуемый размер изображений
img_width = 200
img_height = 50

# Во сколько раз изображение будет уменьшено.
# У нас будет 2 свёрточных блока,
# у каждого блока будет слой пулинга, который субдескритизирует
# изображение в 2 раза, поэтому итоговый коэффициент будет 4.
downsample_factor = 4

# Максимальная возможная длина капчи
max_length = max([len(label) for label in labels])
\end{minted}
\par Затем создаются функции для сопоставления символов индексам и наоборот. Это распространённая практика в сфере обработки текстов, чтобы было проще работать с данными в виде чисел.
\begin{minted}[fontsize=\footnotesize,breaklines]{python}
# Сопоставление символов к числам
char_to_num = tf.keras.layers.StringLookup(vocabulary=list(characters), mask_token=None)

# Сопоставление чисел обратно к символам
num_to_char = tf.keras.layers.StringLookup(
    vocabulary=char_to_num.get_vocabulary(), mask_token=None, invert=True
)
\end{minted}
\par После этого создаются два набора данных - для обучения модели и для валидации результатов. Считанные изображения перемешиваются и распределяются поровну в два набора.
\begin{minted}[fontsize=\footnotesize,breaklines]{python}
def split_data(images, labels, train_size=0.8, test_size=0.1, shuffle=True):
    # 1. Берём количество всех картинок
    size = len(images)
    # 2. Делаем массив индексов и, если надо, перемешиваем его
    indices = tf.keras.ops.arange(size)
    if shuffle:
        tf.keras.random.shuffle(indices)
    # 3. Берём размер тренировочной выборки
    train_samples = int(size * train_size)
    test_samples = train_samples + int(size * test_size) + 2
    print(train_samples, test_samples)
    # 4. Разделяем данные в наборы для обучения и валидации
    x_train, y_train = images[indices[:train_samples]], labels[indices[:train_samples]]
    x_valid, y_valid = images[indices[train_samples:test_samples]], labels[indices[train_samples:test_samples]]
    x_test, y_test = images[indices[test_samples:]], labels[indices[test_samples:]]
    return x_train, x_valid, y_train, y_valid, x_test, y_test


x_train, x_valid, y_train, y_valid, x_test, y_test = split_data(np.array(images), np.array(labels))
print(len(x_train), len(x_valid), len(y_train), len(y_valid), len(x_test), len(y_test))
\end{minted}
\par Далее создаётся функция непосредственно считывания изображения и приведения его в формат, удобный для нейронной сети. Изображение конвертируется в монохромный формат, превращается в матрицу со значениями от 0 до 1, и транспонируется, чтобы модель считывала изображения вдоль ширины, слева направо, вместо того чтобы считывать всё изображение разом.
\begin{minted}[fontsize=\footnotesize,breaklines]{python}
def encode_single_sample(img_path, label):
    # 1. Считываем изображение
    img = tf.io.read_file(img_path)
    # 2. Декодируем и конвертируем в монохром
    img = tf.io.decode_png(img, channels=1)
    # 3. Конвертируем в дробь от 0 до 1
    img = tf.image.convert_image_dtype(img, tf.float32)
    # 4. Изменяем размер изображений под требуемый
    img = tf.keras.ops.image.resize(img, [img_height, img_width])
    # 5. Транспонируем изображение, чтобы ось времени
    # соответствовала ширине изображения
    img = tf.keras.ops.transpose(img, axes=[1, 0, 2])
    # 6. Сопоставляем символы к числам
    label = char_to_num(tf.strings.unicode_split(label, input_encoding="UTF-8"))
    # 7. Возвращаем словарь с двумя вводными данными для модели
    return {"image": img, "label": label}
\end{minted}
\par Затем создаются два объекта наборов данных TensorFlow из тех наборов данных, которые были приготовлены ранее.
\begin{minted}[fontsize=\footnotesize,breaklines]{python}
train_dataset = tf.data.Dataset.from_tensor_slices((x_train, y_train))
train_dataset = (
    train_dataset.map(encode_single_sample, num_parallel_calls=tf.data.AUTOTUNE)
    .batch(batch_size)
    .prefetch(buffer_size=tf.data.AUTOTUNE)
)

validation_dataset = tf.data.Dataset.from_tensor_slices((x_valid, y_valid))
validation_dataset = (
    validation_dataset.map(encode_single_sample, num_parallel_calls=tf.data.AUTOTUNE)
    .batch(batch_size)
    .prefetch(buffer_size=tf.data.AUTOTUNE)
)
\end{minted}
\item Создание модели
\par Сперва создаётся слой вводных данных, содержащий изображения Captcha и текст ответов на них.
\par Затем добавляются два двухмерных свёрточных слоя, с количеством выходных фильтров, равных 32, размером свёрточного окна 2 на 2, с линейным выпрямителем (Rectified linear unit, ReLU) в качестве функции активации и инициализатором хи-нормаль.
\par Затем добавляются слои пулинга, для субдескритизации изображения в 4 раза.
\par После этого добавляются два двунаправленных рекуррентных слоя, использующих архитектуру долгой краткосрочной памяти (Long short-term memory, LSTM). Они проходят по изображению слева направо, используя ширину в качестве единицы времени.
\par Затем добавляется полносвязный слой для агрегации выходных данных рекурретных слоёв.
\par Далее добавляется слой CTC, оценивающий потерю.
\begin{minted}[fontsize=\footnotesize,breaklines]{python}
def build_model():
    # Ввод данных в модель
    input_img = tf.keras.layers.Input(
        shape=(img_width, img_height, 1), name="image", dtype="float32"
    )
    labels = tf.keras.layers.Input(name="label", shape=(None,), dtype="float32")

    # Первый свёрточный слой
    x = tf.keras.layers.Conv2D(
        32,
        (3, 3),
        activation="relu",
        kernel_initializer="he_normal",
        padding="same",
        name="Conv1",
    )(input_img)
    x = tf.keras.layers.MaxPooling2D((2, 2), name="pool1")(x)

    # Второй свёрточный слой
    x = tf.keras.layers.Conv2D(
        64,
        (3, 3),
        activation="relu",
        kernel_initializer="he_normal",
        padding="same",
        name="Conv2",
    )(x)
    x = tf.keras.layers.MaxPooling2D((2, 2), name="pool2")(x)

    # Пулинг с размером фильтра и шага 2, с функцией максимума.
    # В итоге, изображение будет в 4 раза меньше.
    new_shape = ((img_width // 4), (img_height // 4) * 64)
    x = tf.keras.layers.Reshape(target_shape=new_shape, name="reshape")(x)
    x = tf.keras.layers.Dense(64, activation="relu", name="dense1")(x)
    x = tf.keras.layers.Dropout(0.2)(x)

    # Рекуррентные слои
    x = tf.keras.layers.Bidirectional(tf.keras.layers.LSTM(128, return_sequences=True, dropout=0.25))(x)
    x = tf.keras.layers.Bidirectional(tf.keras.layers.LSTM(64, return_sequences=True, dropout=0.25))(x)

    # Слой вывода
    x = tf.keras.layers.Dense(
        len(char_to_num.get_vocabulary()) + 1, activation="softmax", name="dense2"
    )(x)

    # Слой классификации по рейтингу, для вычисления потерь при каждом шаге
    output = CTCLayer(name="ctc_loss")(labels, x)

    # Задаём модель
    model = tf.keras.models.Model(
        inputs=[input_img, labels], outputs=output, name="ocr_model_v1"
    )
    # Оптимизатор
    opt = tf.keras.optimizers.Adam()
    # Компилируем модель
    model.compile(optimizer=opt)
    return model
\end{minted}
\item Обучение модели
\par Обучение происходило в 100 этапов, с использованием тренировочного и валидационного наборов, подготовленных ранее. На первых этапах потери были высокими, но постепенно стали близки к нулю, говоря о хороших результатах обучения.
\\
\begin{figure}[H]
  \centering
  \includegraphics[width=0.7\textwidth]{graph.png}
  \caption{График обучения}
\end{figure}
\item Вывод данных
\par Последним этапом было создание функции для расшифровки выходных данных модели и преобразования их в текст, с помощью которого можно будет самостоятельно проверить, насколько хорошо модель справляется с определением текста на заданном изображении.
\begin{minted}[fontsize=\footnotesize,breaklines]{python}
# Расшифровка прогноза
def decode_batch_predictions(pred):
    input_len = np.ones(pred.shape[0]) * pred.shape[1]
    results = ctc_decode(pred, input_length=input_len, greedy=True)[0][0][
        :, :max_length
    ]
    # Пробегаемся по результатам и возвращаем текст
    output_text = []
    for res in results:
        res = tf.strings.reduce_join(num_to_char(res)).numpy().decode("utf-8")
        output_text.append(res)
    return output_text
\end{minted}
\begin{figure}[H]
  \centering
  \includegraphics[width=\textwidth]{results.png}
  \caption{Примеры результатов}
\end{figure}
\end{enumerate}
\end{document}
